\chapter{关系查询、指标计算与语义服务}
\label{ch:week03-sql-kpi}

企业数据模型规定记录的结构，指标查询则把记录汇总为特定时期、产品或部门的
经营统计量。SQL 提供筛选、连接与聚合运算，指标定义规定运算的对象和含义。
比率、快照和窗口统计采用不同的汇总方法；制造指标还需明确数量、时间与对象
范围。自然语言查询须依据这些定义生成 SQL、检验结果，并满足数据库执行的
权限和资源限制。

\section{查询任务与数据模型}

SQL 是关系数据库的声明式查询语言。销售额、延期率和高风险订单数等统计量，均可通过记录筛选、表间连接和分组聚合计算。查询文本保存了这些计算的定义，便于重复执行和核对。

事实表粒度决定记录与计数单位，其定义见第~\ref{ch:week02-enterprise-data-model}~章。若订单表的粒度是一行一个订单，那么订单数可以直接按行计数；若事实表的粒度是一行一个订单明细项，那么订单数需要按订单编号去重；若同一张表混入订单头、订单行和库存快照，就会导致重复计数。SQL 语法本身不会阻止这种错误，只有明确的业务粒度和指标定义才能防止它。

\begin{definition}[指标层]
指标层是组织对管理指标的统一表达层。它规定指标名称、业务定义、对象范围、事实粒度、分子分母、维度、时间窗口、单位、刷新频率、权限和责任人，并提供可复用的 SQL 或语义配置。
\end{definition}

指标查询按以下依赖关系构造。
\[
\text{业务对象}
\rightarrow
\text{事实粒度}
\rightarrow
\text{指标口径}
\rightarrow
\text{查询表达}
\rightarrow
\text{管理解释}
\rightarrow
\text{行动审计}.
\]

管理指标查询的正确性包括语法有效、业务问题匹配、数据范围一致、分子分母明确以及结果可复核。

\section{关系查询与聚合运算}
\subjectindex{SQL}{SQL}
\subsection{SQL 查询的关系语义}
\subjectindex{guanxidaishu}{关系代数}

关系模型把数据看作关系，即带有属性名的元组集合。SQL 查询可以理解为在关系上的变换。尽管实际数据库会优化执行计划，但从逻辑上看，一个典型查询依次经历表来源、连接、过滤、分组、聚合、筛选和输出。

\begin{definition}[SQL 查询的逻辑处理顺序]
一个常见 SQL 查询的逻辑顺序可概括为：
\[
\texttt{FROM/JOIN}
\rightarrow
\texttt{WHERE}
\rightarrow
\texttt{GROUP BY}
\rightarrow
\texttt{HAVING}
\rightarrow
\texttt{SELECT}
\rightarrow
\texttt{ORDER BY}.
\]
这一顺序表示查询语义的理解方式，不一定等于数据库系统实际执行计划。
\end{definition}

例如，下面的查询按区域计算 2026 年 3 月订单收入：

\begin{lstlisting}[style=sqlstyle,caption={区域销售额查询}]
SELECT
  c.region,
  SUM(o.quantity * p.unit_price) AS revenue
FROM orders AS o
JOIN customers AS c
  ON o.customer_id = c.customer_id
JOIN products AS p
  ON o.product_id = p.product_id
WHERE o.order_date >= DATE '2026-03-01'
  AND o.order_date <  DATE '2026-04-01'
GROUP BY c.region
ORDER BY revenue DESC;
\end{lstlisting}

该查询确定了四项计算约定。第一，分析对象是 2026 年 3 月订单；第二，订单通过客户维度获得地区，通过产品维度获得售价；第三，收入按地区汇总；第四，输出排序用于比较区域经营表现。样例数据假定每张订单只有一种产品，且产品单价在考察期内不变，因此可以用数量乘该单价。真实交易收入应优先使用订单行上留存的成交价格，并说明折扣、退货与税费口径；历史销售额应由当时的成交条件确定。

\begin{remark}
\texttt{WHERE} 在聚合前过滤原始事实记录，\texttt{HAVING} 在聚合后过滤分组结果。若把二者混淆，就可能把“筛选订单”和“筛选地区”混为一谈。前者改变事实集合，后者改变展示集合。
\end{remark}

\subsection{关系代数视角}

SQL 建立在关系模型和集合运算之上。关系代数可以描述查询对事实集合的变换，用于检查对象遗漏和重复计数。

\begin{table}[htbp]
\centering
\caption{关系代数操作与 SQL 子句}
\begin{tabular}{P{0.20\linewidth}P{0.29\linewidth}P{0.36\linewidth}}
\toprule
\rowcolor{BDMTableHead}
关系操作 & SQL 表达 & 管理含义\\
\midrule
选择 \(\sigma\) & \texttt{WHERE} & 限定事实集合，例如只保留某月订单。\\
投影 \(\pi\) & \texttt{SELECT} & 选择输出属性，例如地区和收入。\\
连接 \(\Join\) & \texttt{JOIN} & 补充业务对象属性，例如客户地区和产品价格。\\
分组聚合 \(\gamma\) & \texttt{GROUP BY} & 将事实提升到分析粒度，例如区域级 KPI。\\
排序 \(\tau\) & \texttt{ORDER BY} & 为比较和展示建立顺序。\\
\bottomrule
\end{tabular}
\end{table}

前一节的区域收入查询可以抽象为
\[
\gamma_{\text{region};\, \sum(q\cdot price)\rightarrow revenue}
\left(
\sigma_{\text{March 2026}}
(\text{orders}\Join \text{customers}\Join \text{products})
\right).
\]
表达式由内向外依次补充维度、选择日期范围并汇总收入。各运算的输入与输出均有确定含义，因而可以分别核对连接后的记录数、筛选范围和分组金额。

\begin{convention}[声明式查询与可审计性]
声明式 SQL 明确事实来源、过滤条件、连接键、分组维度和聚合方式，查询文本因此构成指标审计的依据。
\end{convention}

\subsection{查询的选择与分组条件}
关系代数把查询视为对集合的运算；实际 SQL 默认保留重复行，更接近多重集语义。
这个差别具有业务后果。两个订单行即使产品、数量和金额相同，也可能是两个不同事实；
对全部字段使用 \texttt{DISTINCT} 并不能可靠地替代业务主键去重。
重复究竟是合法重复还是采集错误，应由第~\ref{ch:week02-enterprise-data-model}~章
定义的对象身份决定。

聚合查询的核查包括四项：候选记录的范围、连接后的事实匹配次数、分组所确定的
结果粒度，以及每组保留的统计状态。数据库执行器可能改变实际计算顺序以提高效率，但正确的优化
必须保持结果语义及指标的分母定义不变。

空值引入三值逻辑。SQL 中 \texttt{NULL = NULL} 的结果为未知；
\texttt{WHERE} 只保留条件为真的行。因此，条件 \texttt{actual\_date > due\_date}
不会自动选出已经到期但尚未交付的订单。把未交付订单排除之后，查询回答的已不再是
“所有到期承诺中有多少延期”，而只覆盖实际交付日期已记录的订单。

\subsection{事实粒度与分析粒度}

指标计算需分别声明源记录、分析结果与展示结果的粒度。

\begin{definition}[事务粒度、分析粒度与展示粒度]
事务粒度是源事实表中一行记录代表的业务事实；分析粒度是 SQL 查询分组后每一行结果代表的业务层级；展示粒度是看板、报告或管理会议中最终呈现的层级。
\end{definition}

例如，订单明细事实表的事务粒度可能是一行一个订单明细项；按地区、月份聚合后的分析粒度是一行一个地区月份；在管理层月报中展示的粒度可能是一行一个区域。三者可以不同，但必须被明确说明。

\begin{table}[htbp]
\centering
\caption{三类粒度的例子}
\begin{tabular}{P{0.21\linewidth}P{0.31\linewidth}P{0.34\linewidth}}
\toprule
\rowcolor{BDMTableHead}
粒度类型 & 示例 & 典型风险\\
\midrule
事务粒度 & 一行一个订单明细项。 & 误把明细行数当订单数。\\
分析粒度 & 一行一个地区、月份。 & 聚合维度遗漏导致混合口径。\\
展示粒度 & 一行一个区域或业务线。 & 从低层百分比直接平均得到高层百分比。\\
\bottomrule
\end{tabular}
\end{table}

构造查询之前，应先确定源事实表和目标结果表的粒度。随后选择维度连接、过滤条件与聚合函数，使每一行输出恰好对应目标分析单位。

\subsection{连接、粒度与重复计数}
\subjectindex{lianjie}{连接}
\subjectindex{chongfujishu}{重复计数}

内连接保留两张表中能够匹配的记录，左连接保留左表全部记录并补充右表信息，外连接可用于识别匹配与未匹配的范围。连接可能复制或遗漏事实记录，因此需要核查其对事实表粒度的影响。

假设订单表 \(\texttt{orders}\) 的粒度是一行一个订单，产品表 \(\texttt{products}\) 的粒度是一行一个产品。如果 \(\texttt{products.product\_id}\) 唯一，且每个订单具有非空、可匹配的产品编号，则按产品编号内连接满足
\[
|\texttt{orders}\Join\texttt{products}|=|\texttt{orders}|.
\]
这说明连接只是补充产品属性，没有复制或遗漏订单事实。唯一性本身只能保证每张订单至多匹配一条产品记录，不能保证一定匹配；存在空外键或孤立外键时，内连接仍会丢行。若某个被订单引用的产品编号出现两次，相应订单则会被复制，销售额、订单数和延期率都可能被扭曲。

\begin{table}[htbp]
\centering
\caption{连接错误与指标后果}
\begin{tabular}{P{0.21\linewidth}P{0.33\linewidth}P{0.34\linewidth}}
\toprule
\rowcolor{BDMTableHead}
问题 & 数据原因 & 指标后果\\
\midrule

事实行膨胀 & 维度键不唯一，或多对多关系未预聚合。 & 销售额、订单量、客户数被重复计算。\\
事实行丢失 & 使用内连接，但维度缺失或外键无法匹配。 & 低估业务规模，遗漏异常对象。\\
时间错配 & 使用当前维度属性解释历史事实。 & 历史地区、等级、价格口径被改写。\\
粒度混合 & 订单头、订单行、库存快照混在同一聚合层。 & KPI 分母不清，比较失真。\\
\bottomrule
\end{tabular}
\end{table}

\begin{proposition}[连接前后的行数检查]
若 \(F\) 是事实表，\(D\) 是维度表，且 \(D\) 在连接键上唯一并覆盖 \(F\) 的外键，则 \(F\) 与 \(D\) 的内连接不会改变事实行数。若连接后行数增加，通常说明维度键不唯一；若连接后行数减少，通常说明存在无法匹配的外键。
\end{proposition}

管理指标查询的连接复核需确认事实表的记录单位、连接后的事实覆盖及维度键的唯一性。

\section{指标的代数结构与制造口径}
\subsection{指标定义的形式化表示}

关键绩效指标（key performance indicator, KPI）的定义包括对象范围、计算方法及使用条件。为便于核对，记
\[
M=(O,\mathcal G,\mathcal T,N,D,u,\rho),
\]
其中 \(O\) 是被度量的对象集合或事件集合，\(\mathcal G\) 是分组维度，\(\mathcal T\) 是时间窗口，\(N\) 是分子或可加度量，\(D\) 是分母，\(u\) 是单位，\(\rho\) 是业务责任或解释责任。

以配套零售样例的“已交付订单准时率”为例。对象集合 \(O\) 是已交付订单，时间窗口 \(\mathcal T\) 必须明确采用订单日期或承诺交付日期，分子 \(N\) 是其中准时交付的订单数，分母 \(D\) 是该范围内已交付订单数，单位 \(u\) 是百分比，责任 \(\rho\) 可能涉及供应链、仓储和物流团队。后文的制造 OTD 则采用到期承诺口径，二者不作为同一指标混用。

\begin{table}[htbp]
\centering
\caption{样例 KPI 字典}
\begin{tabular}{P{0.18\linewidth}P{0.25\linewidth}P{0.39\linewidth}}
\toprule
\rowcolor{BDMTableHead}
指标 & 业务含义 & 计算口径\\
\midrule
销售额 & 订单带来的标准收入。 & \(\sum q_i \cdot price_i\)\\
毛利 & 收入扣除标准成本后的贡献。 & \(\sum q_i(price_i-cost_i)\)\\
毛利率 & 收入中转化为毛利的比例。 & \(\sum GP_i / \sum Revenue_i\)\\
已交付订单准时率 & 已交付订单中按承诺日期交付的比例。 & \(\sum I(a_i\le p_i)/n\)\\
平均延期天数 & 订单超过承诺日期的平均延迟。 & \(\sum \max(a_i-p_i,0)/n\)\\
\bottomrule
\end{tabular}
\end{table}

毛利率和准时交付率是比率型指标。跨区域、产品或月份汇总时，应分别合并分子、分母，再在目标粒度上计算比率，
\[
\frac{\sum_{g}N_g}{\sum_{g}D_g}
\quad\text{通常不等于}\quad
\frac{1}{|\mathcal G|}\sum_g\frac{N_g}{D_g}.
\]

\begin{remark}
左边是按事实量加权后的总体比率，右边是分组百分比的平均数。除非每个分组分母相同，否则二者通常不同。
\end{remark}

\subsection{可合并的统计状态}
\subjectindex{kehebingzhuangtai}{可合并统计状态}
一个指标在日、周、工站和区域之间汇总时，需要保留足够的中间统计量。设不相交事实集合 \(D_1,D_2\) 的汇总状态为 \(S(D)\)。若存在结合运算 \(\oplus\)，满足
\[
S(D_1\cup D_2)=S(D_1)\oplus S(D_2),
\]
便能先局部计算再合并。总量的状态是和；均值的状态是 \((\sum x,n)\)；比率的状态是 \((N,D)\)。最终展示函数分别为恒等、\(s/n\) 与 \(N/D\)。平均值自身通常不满足上述合并律。

\begin{example}[良率的正确汇总]
甲线产出 100 件、良品 90 件，乙线产出 900 件、良品 720 件。平均两条线的良率得到 \(85\%\)，整体良率却是
\[
\frac{90+720}{100+900}=81\%.
\]
两条线平均值回答的是“随机选一条线的平均良率”，整体比率回答的是“随机选一件产品为良品的概率”。前者以产线为抽样单位，后者以产品件数为抽样单位。
\end{example}

连接也需要代数审计。若事实 \(i\) 的金额为 \(x_i\)，内连接后匹配到 \(m_i\) 条记录，直接求和得到 \(\sum_i m_ix_i\)，未必等于 \(\sum_i x_i\)。维度唯一性与完整匹配共同保证 \(m_i=1\)；订单连接多条质检事件则需要先按订单聚合事件，或改用半连接筛选。随意使用 \texttt{DISTINCT} 可能恰好掩盖样例错误，却会误删金额相同的不同订单。

\subsection{制造指标的计算}
\subjectindex{Yield}{Yield（良率）}
\subjectindex{OEE}{OEE（设备综合效率）}
\subjectindex{WIP}{WIP（在制品）}
\subjectindex{OTD}{OTD（准时交付率）}
\subjectindex{CycleTime}{Cycle Time（周期）}
\subjectindex{lianglv}{良率}
\subjectindex{shebeizonghexiaolv}{设备综合效率}
\subjectindex{zaizhipin}{在制品}
\subjectindex{zhunshijiaofulv}{准时交付率}
\subjectindex{zhouqi}{周期}
约定观察窗口为 \([t_0,t_1)\)，所有数量使用相同产品计量单位。先计算可相加的事实，再生成比率：
\[
\mathrm{Yield}=\frac{\sum_i G_i}{\sum_i N_i},\qquad
\mathrm{CT}=\frac{\sum_i n_i\bar c_i}{\sum_i n_i}.
\]
这里 \(G_i,N_i\) 是同一检验范围的良品与总件数；\(n_i\) 是完成批次数，\(\bar c_i\) 是这些批次的平均周期。若只知道各组均值而不知道样本量，就无法还原整体均值。Cycle Time 只观察完成批次还可能遗漏长期滞留批次，因此应并列报告未完成批次的年龄。

对于单一节拍、同一产出范围的设备，设计划时间 \(T_p\)、运行时间 \(T_r\)、理想单件时间 \(t_c\)，则
\[
A=\frac{T_r}{T_p},\quad
P=\frac{Nt_c}{T_r},\quad
Q=\frac GN,\quad
\mathrm{OEE}=APQ=\frac{Gt_c}{T_p}.
\]
时间单位必须统一；理想节拍乘件数必须与运行时间比较。对多产品不同节拍的情形，一个可解释的聚合是
\[
\mathrm{OEE}_{\rm time}=\frac{\sum_i t_{c,i}G_i}{\sum_i T_{p,i}}.
\]
若坚持 \(A P Q\) 分解，质量项应使用理想工时加权的 \(Q=\sum_i t_{c,i}G_i/\sum_i t_{c,i}N_i\)，不能将按件数汇总的良率无条件乘入。上述等式还要求产品之间的计划时间正确分摊，没有重复占用设备日历\cite{semi03}。

WIP 是时点存量，其时间平均应定义为
\[
\overline{\mathrm{WIP}}
=\frac1{t_1-t_0}\int_{t_0}^{t_1}\mathrm{WIP}(t)\,\mathrm dt.
\]
等间隔且完整的快照可用算术平均近似；间隔不等时应按持续时间加权。若每个工站每天包含多个产品族，应先在日内对互斥产品族求和，再跨天平均。直接平均所有行得到的是产品族快照的平均值；工站总在制品则由同日各产品族数量求和得到。

采用到期承诺口径时，OTD 的分母是窗口内到期且属于统计范围的订单数。到统计截止日仍未完成的逾期订单属于未按时履约；尚未到期的订单不属于同一风险集合。若使用“已交付订单准时率”，应另命名并解释其可能排除长期拖欠订单的选择偏差。

该查询采用一行一个有效订单承诺的关系模式，\texttt{delivered\_at} 表示截至统计时点已知的实际交付时间。查询前须按既定规则选择承诺版本并保留变更记录，以便区分服务改善与承诺日期调整。

\subsection{条件聚合与安全除法}

SQL 中的 KPI 通常通过条件聚合表达。下面的查询计算区域级销售额、毛利、准时交付率与平均延期天数。

\begin{lstlisting}[style=sqlstyle,caption={区域级运营 KPI 查询}]
WITH order_enriched AS (
  SELECT
    o.order_id,
    c.region,
    o.quantity * p.unit_price AS revenue,
    o.quantity * (p.unit_price - p.unit_cost) AS gross_profit,
    CASE
      WHEN o.actual_date <= o.promised_date THEN 1
      ELSE 0
    END AS on_time_flag,
    GREATEST(date_diff('day', o.promised_date, o.actual_date), 0) AS delay_days
  FROM orders AS o
  JOIN customers AS c
    ON o.customer_id = c.customer_id
  JOIN products AS p
    ON o.product_id = p.product_id
  WHERE o.actual_date IS NOT NULL
)
SELECT
  region,
  COUNT(*) AS delivered_orders,
  SUM(revenue) AS revenue,
  SUM(gross_profit) AS gross_profit,
  SUM(gross_profit) / NULLIF(SUM(revenue), 0) AS gross_margin,
  SUM(on_time_flag) * 1.0 / NULLIF(COUNT(*), 0) AS on_time_rate,
  AVG(delay_days) AS average_delay_days
FROM order_enriched
GROUP BY region
ORDER BY revenue DESC;
\end{lstlisting}

\texttt{NULLIF}(x,0) 在 \(x=0\) 时返回空值，其余情况返回 \(x\)。用它处理分母，可将零分母比率报告为未定义。例如，某区域没有已交付订单时，准时率未定义；有已交付订单但准时数为0时，准时率才为0。

\begin{definition}[安全除法]
安全除法是指在计算比率指标时显式处理分母为零、分母为空或分子分母口径不一致的情况。常见做法是返回空值、单独标记或在指标字典中说明不可定义状态。
\end{definition}

\subsection{空值、未完成事件与业务状态}
\subjectindex{kongzhi}{空值}

SQL 中的 \texttt{NULL} 不等同于数值 0，也不等同于空字符串。它通常表示未知、缺失、尚未发生或不适用。在管理指标中，空值处理必须回到业务状态，而不能只按技术便利处理。

\begin{table}[htbp]
\centering
\caption{空值与业务解释}
\begin{tabular}{P{0.24\linewidth}P{0.28\linewidth}P{0.34\linewidth}}
\toprule
\rowcolor{BDMTableHead}
情形 & 常见 SQL 处理 & 管理解释\\
\midrule
未交付订单无实际日期 & 已交付口径排除；到期口径保留已到期订单。 & 已到期且确认未交付是履约失败；尚未到期或回写缺失需分别判断。\\
区域没有已交付订单 & 比率返回 \texttt{NULL}。 & 指标未定义，单独报告业务数量为0。\\
产品成本缺失 & 标记数据质量问题。 & 毛利无法可信计算，需要补数或剔除说明。\\
客户维度缺失 & 左连接保留事实并标为未知客户。 & 业务规模存在，但归属和分层解释受限。\\
\bottomrule
\end{tabular}
\end{table}

\begin{remark}
一个严谨的指标系统应区分三类状态：“业务没有发生”“业务发生但结果失败”“数据缺失导致无法判断”。这三者在管理上对应完全不同的行动。
\end{remark}

\subsection{一个到期队列的完整计算}
\label{sec:due-cohort-example}
\subjectindex{daoqiduilie}{到期队列}
设统计时点为9月8日结束，统计对象是9月1日至7日到期的四张订单。其状态如下；
日期、订单和数值均为假设值，承诺日期采用统计时点前已经批准的版本。
\begin{table}[htbp]
\centering
\caption{到期队列与已交付队列回答不同问题}
\label{tab:due-cohort}
\begin{tabular}{llll}
\toprule
\rowcolor{BDMTableHead}
订单 & 承诺日期 & 实际交付日期 & 截止时点判断\\
\midrule
A & 9月2日 & 9月2日 & 准时\\
B & 9月4日 & 9月6日 & 延期但已交付\\
C & 9月6日 & 尚未交付 & 已到期未交付\\
D & 9月7日 & 9月7日 & 准时\\
\bottomrule
\end{tabular}
\end{table}

按到期承诺评价，分母为4，分子为2，OTD 为50\%。若只对已交付的 A、B、D 计算，
分母为3，得到约66.7\%。两个数都可以是各自定义下的正确计算，却不能用于同一个
不加说明的“交付表现”结论。若 C 在9月10日交付，到期队列中的历史承诺完成情况
仍不能因此从延期变成准时；它改变的是最终完成日期和延期天数。

只保存50\%无法判断样本量，也无法检查迟到数据；同时保存到期订单数、准时数、
未交付数和数据截止时点，可以核对后续重算的变化。尚未完成的对象仍属于相应
统计范围，其当前状态与已完成结果需要分别记录。缺失选择的影响见
第~\ref{ch:week04-data-quality}~章。

\section{时间窗口与指标核查}
\subsection{窗口函数与经营节奏}
\subjectindex{chuangkohanshu}{窗口函数}

聚合查询将同组多行压缩为一行，窗口函数在保留行身份的同时，对与当前行相关的一组行进行计算。经营分析由此可以同时保留单个订单及其在客户历史、区域排名或时间趋势中的位置。

\begin{lstlisting}[style=sqlstyle,caption={客户最近三笔订单收入}]
SELECT
  o.customer_id,
  o.order_id,
  o.order_date,
  o.quantity * p.unit_price AS revenue,
  SUM(o.quantity * p.unit_price) OVER (
    PARTITION BY o.customer_id
    ORDER BY o.order_date, o.order_id
    ROWS BETWEEN 2 PRECEDING AND CURRENT ROW
  ) AS customer_recent_revenue,
  ROW_NUMBER() OVER (
    PARTITION BY o.customer_id
    ORDER BY o.order_date, o.order_id
  ) AS customer_order_sequence
FROM orders AS o
JOIN products AS p
  ON o.product_id = p.product_id
ORDER BY o.customer_id, o.order_date, o.order_id;
\end{lstlisting}

窗口函数在管理分析中常用于：
\begin{itemize}
  \item 对客户、产品、供应商或区域进行排名；
  \item 计算累计销售额、累计延期订单数和累计缺货次数；
  \item 计算移动窗口收入、近三次订单金额或近七天投诉数；
  \item 识别同一对象的前后状态变化，例如订单状态流转或库存变化。
\end{itemize}

窗口函数的风险在于窗口定义不清。若 \texttt{ORDER BY} 没有唯一排序键，排名和前后项可能不稳定；若窗口框架没有明确说明行数或时间范围，累计值可能被误解为移动值；若分区维度错误，客户级行为会被混成全局行为。

\subsection{指标计算的核查}

指标层集中保存指标定义及其计算实现。每个指标应记录定义负责人、数据源、事实粒度、分析维度、可加性和测试结果。引用指标的团队依据同一版本的定义和数据范围复现结果，并据此核查差异。

一个基础指标配置可以写成如下结构：

\begin{lstlisting}[style=bookcode,language={},caption={指标配置示例}]
metrics:
  - name: gross_margin
    label: Gross margin
    type: ratio
    numerator: gross_profit
    denominator: revenue
    grain: order
    time_dimension: order_date
    dimensions: [region, segment, category, channel]
    owner: finance
\end{lstlisting}

\begin{table}[htbp]
\centering
\caption{指标定义的基本字段}
\begin{tabular}{P{0.22\linewidth}P{0.60\linewidth}}
\toprule
\rowcolor{BDMTableHead}
字段 & 说明\\
\midrule
指标名称 & 组织中统一使用的名称。\\
业务定义 & 指标回答的管理问题。\\
事实粒度 & 计算依据的一行事实代表什么。\\
分子与分母 & 比率指标必须保存分子和分母定义。\\
时间窗口 & 使用订单日期、交付日期、快照日期还是入库日期。\\
维度范围 & 可以按哪些对象、组织层级和时间层级下钻。\\
责任人与版本 & 谁维护口径，何时修改，如何复核。\\
\bottomrule
\end{tabular}
\end{table}

指标层使模型、看板和报告共用已确认的定义，减少重复计算中的口径分歧，并为复核提供版本依据。

\subsection{从指标结果到管理解释}

指标结果的管理解释应明确事实范围、比较基准、计算口径、异常位置、可行行动及复核依据。

\begin{table}[htbp]
\centering
\caption{KPI 结果解释模板}
\begin{tabular}{P{0.20\linewidth}P{0.63\linewidth}}
\toprule
\rowcolor{BDMTableHead}
解释环节 & 需要回答的问题\\
\midrule
事实描述 & 指标在什么对象、时间窗口和维度上发生了什么变化？\\
比较基准 & 与上期、目标值、预算、同类区域或历史分位数相比如何？\\
口径说明 & 分子、分母、排除条件、空值和未完成事件如何处理？\\
异常定位 & 哪些客户、产品、渠道、供应商或流程节点贡献最大？\\
行动建议 & 可以采取哪些行动，行动对象和约束是什么？\\
复核路径 & SQL、输入表版本、指标字典和责任人能否被追溯？\\
\bottomrule
\end{tabular}
\end{table}

例如，分析南区准时交付率时，可先报告其与全体水平的差异，再按订单优先级和产品类别分解。若使用已交付订单口径，应同时列出尚未交付订单数。随后结合供应商交期、出库时间和运输记录检验各项解释，并识别需要进一步处理的订单。

\section{自然语言查询与 SQL 生成}

\subsection{自然语言问题的形式化}
\subjectindex{NL2SQL}{NL2SQL}
\subjectindex{AIwenshu}{AI问数}
“找出最近良率最差的工站”至少留下四个未定项：最近是自然周还是滚动七天，良率是首检还是最终良率，最差按水平还是相对目标偏差，低产量工站是否有最低分母要求。语言模型可以生成流畅 SQL，但没有足够信息时，应返回待澄清项或明示默认口径。

Spider 2.0 研究企业级数据和多步查询任务，BIRD-INTERACT 研究动态交互条件下的问数任务\cite{spider03,bird03}。这些工作将任务定义、模式探索和执行反馈纳入系统评价。企业应用可据自身数据模式和用户问题构造相应测试集。

\begin{definition}[查询请求]
将查询请求记为
\[
q=(m,g,w,f,v,u,\mathcal E),
\]
分别表示指标、分组粒度、时间窗口、过滤条件、指标定义版本、用户身份和预期输出条件。查询请求规定计算目标，SQL 给出该目标的具体实现。
\end{definition}

\begin{figure}[htbp]
\centering
\begin{tikzpicture}[node distance=0.65cm]
\node[bdmnode] (q) {查询请求\\指标与时间};
\node[bdmnode,right=of q] (s) {语义检索\\表与字段};
\node[bdmnode,right=of s] (c) {候选 SQL\\静态检查};
\node[bdmnode,below=of c] (e) {受限执行\\预算与权限};
\node[bdmnode,left=of e] (v) {结果验证\\口径与不变量};
\node[bdmnode,left=of v] (r) {证据回答\\数值与来源};
\draw[bdmarrow] (q)--(s);
\draw[bdmarrow] (s)--(c);
\draw[bdmarrow] (c)--(e);
\draw[bdmarrow] (e)--(v);
\draw[bdmarrow] (v)--(r);
\end{tikzpicture}
\caption{自然语言查询的生成、执行与结果检验}
\end{figure}

\subsection{制造查询的验证步骤}
问题为“上周哪个工站最终良率最低”。第一步固定自然周起止、时区、最终检验范围和最小件数。第二步查找已确认的指标定义及实际字段。第三步生成按 \texttt{work\_center} 汇总良品和总数的 SQL。第四步核对零分母、负数、重复行和产量筛选。第五步返回工站、分子、分母、比率、数据截止时间与查询版本。

结果还需要不变量测试：良品不能超过同范围总数；工站分组的件数之和应等于未分组件数；将同一个输入文件按行拆成两部分，合并统计状态后应与全量计算一致。改变输入顺序不应改变非顺序指标。此类变形测试可以发现单个标准答案未覆盖的错误。

若进一步研究设备维修，需要把工站良率与批次、缺陷和设备日志关联，确定待检验的故障机制，再估计维修效果与成本。工站级汇总提供异常位置，设备级分析提供维修决策所需的资料。

\subsection{查询执行的权限控制}
\subjectindex{zhixinganquan}{执行安全}
\subjectindex{shaxiang}{沙箱}
查询权限需要同时约束语句、函数和外部资源访问。分析数据库的查询可能调用外部文件读取、网络访问或扩展函数。DuckDB 官方安全文档明确讨论了外部访问与扩展的能力边界\cite{duckdb03}。查询服务宜采用由服务端维护的参数化查询模板、只读数据快照、许可表、行数和时间预算，以及受限文件系统与网络环境。语句解析、数据库权限和运行环境隔离应同时生效。

参数绑定适用于值，表名与列名应从白名单映射；数据库权限必须实际生效；执行日志保存用户问题、规范化查询请求、SQL、参数、数据版本和结果摘要。模型生成的代码经过验证后执行，检索内容作为查询资料使用；系统权限由独立的授权配置决定。

\section{语义查询与企业数据服务}

\subsection{Snowflake 与 Databricks 的产品机制}
\subjectindex{Snowflake}{Snowflake}
\subjectindex{Databricks}{Databricks}
\subjectindex{biaozhunchaxun}{标准查询}
Snowflake Cortex Analyst 的 Verified Query Repository 将经审核的自然语言问题与 SQL 配对，作为可复用的语义示例\cite{snowflake03}。Databricks Genie 的评测文档使用问题集及同义改写检查回答质量\cite{genie03}。前者保存查询示例，后者检验不同提问方式下的结果一致性。

“周良率”“到期 OTD”“日均 WIP”的查询验证可分别保存请求定义、SQL 和小型反例数据。测试时使用未见过的同义问法与不同数据规模，避免只记住固定答案。标准查询库可以提高稳定性；一旦口径或模式变化，它也需要重新验证。

\subsection{本地部署的数据与模型服务}
\subjectindex{bendihuafuwu}{本地化服务}

本地部署需分别确定数据存储、检索计算、模型推理及日志遥测的位置。一种部署方式是在组织内部保存数据库，由语义服务提供指标模板，本地推理服务生成查询请求，执行服务完成授权和计算，再将所需聚合结果用于解释。

Oracle 的 AI Database 26ai 文档和 Private AI Services Container 介绍了向量处理及企业环境中的 AI 服务部署方式\cite{oracle03,oracle10}。部署分析需同时考虑数据位置、推理位置、网络通信和运行日志的保存位置。

若模型有 \(P\) 个参数，每个权重使用 \(b\) 位，纯权重存储约为 \(Pb/8\) 字节。运行还需键值缓存（KV cache）、中间计算和批处理内存。部署评价应在给定上下文长度及并发量下测量正确率、延迟、资源消耗和错误损失。

\subsection{多平台问数系统的可比较任务}
\subjectindex{yuyituxiao}{语义视图}
Databricks Genie Agents 的公开文档要求分析者选定 Unity Catalog 中的数据集，
并提供业务表达式、示例 SQL 与术语说明\cite{databricksGenie2026}。Snowflake
Cortex Analyst 依赖 Semantic Views 中的关系、事实和指标定义，生成 SQL 后
由数据库执行\cite{snowflakeSemantic2026}。两种实现都把业务定义纳入查询
上下文，具体的数据集、权限和计算接口仍由各自平台管理。Snowflake Cortex
Agents 还可把结构化数据查询与文档检索放在同一任务中；一次回答应分别标记
数值来源与文档来源\cite{snowflakeAgents2026}。

企业比较问数产品时，可固定一组问题 \(q_1,\ldots,q_n\) 与独立核对的答案
\(y_1^*,\ldots,y_n^*\)。每题保留数据库快照、语义定义版本、用户角色、
预计可访问的对象及可接受的数值误差。设 \(C_i=1\) 表示第 \(i\) 题的数值和
口径均正确，\(P_i=1\) 表示没有超出授权范围，\(D_i=1\) 表示在期限内返回，
则符合全部要求的任务比例为
\[
  \widehat p_{\rm valid}=\frac1n\sum_{i=1}^{n} C_iP_iD_i.
\]
评价集要包含空分母、迟到记录、同一订单拆成多个批次和用户无权查看的客户。
错误查询即使得到合理数字，也应按 \(C_i=0\) 记录；拒绝越权请求可能是正确的
安全结果，应与可回答任务分开汇报。问题族及数据库规模不同会改变结果，故报告
须列明样本构成，并分别记录准确率、拒绝率、端到端时延与单位任务费用。

\section{订单指标的计算实例}
\subsection{订单指标实例}

订单指标的样例数据位于 \texttt{data/sample/}。表结构如下：
\begin{itemize}
  \item \path{data/sample/customers.csv}：客户维度，一行一个客户；
  \item \path{data/sample/products.csv}：产品维度，一行一个产品，包含标准售价和单位成本；
  \item \path{data/sample/orders.csv}：订单事实，一行一个订单；
  \item \path{data/sample/inventory.csv}：库存快照，一行一个产品、仓库、日期。
\end{itemize}

计算实验分为四步：
\begin{enumerate}
  \item 加载 CSV 数据并检查行数、主键唯一性和外键匹配；
  \item 用 SQL 连接订单、客户和产品三张表，构造订单级收入、毛利和交付状态；
  \item 按区域、客户分层、产品类别和渠道计算 KPI；
  \item 将 KPI 口径写入 \path{src/bdm_decision/warehouse/metrics.yml}，并用 Python 复核关键指标。
\end{enumerate}

样例订单表保存数量，产品表保存考察期内固定的售价与成本。连接后计算销售额与毛利，再按地区或客户分层聚合；历史价格变化的情形采用订单行留存价格。

\subsection{配套程序}

制造指标程序位于 \path{src/bdm_decision/cases/week03_semiconductor_kpi.py}。复现时须按函数定义区分按记录平均的 WIP、按记录数确定的滚动窗口，以及使用件数良率分解的 OEE。工站日总 WIP 需要先合并同日产品族，异质节拍下的 OEE 使用前文的理想工时权重。零分母表示指标未定义，应在报告中列明；程序中的零值返回约定见附录~\ref{app:reproduction}。

订单指标的 Python 复核代码位于：
\begin{center}
\path{src/bdm_decision/cases/week03_sql_kpi.py}
\end{center}

核心函数包括：
\begin{itemize}
  \item \texttt{enrich\_orders()}：连接订单、客户和产品，生成订单级收入、毛利和交付状态；
  \item \texttt{summarize\_by()}：在指定分组维度上计算订单数、收入、毛利率、准时交付率和延期天数；
  \item \texttt{rolling\_customer\_revenue()}：复现窗口函数思想，计算客户最近三笔订单收入。
\end{itemize}

配套 SQL 位于：
\begin{center}
\texttt{src/bdm\_decision/warehouse/build\_kpi\_tables.sql}
\end{center}

该脚本以 DuckDB 风格读取 CSV，建立 \texttt{order\_enriched}、\texttt{kpi\_by\_region}、\texttt{kpi\_by\_month} 和 \texttt{customer\_order\_window} 等视图，用于复现前述 SQL 查询。

\subsection{查询结果的核查}

查询复核包括数据范围、关系约束和聚合结果三个方面，具体检查如下。
\begin{enumerate}
  \item 每个查询说明事实表粒度；
  \item 每个连接说明连接键及其唯一性假设；
  \item 每个比率指标保留分子和分母；
  \item 每个时间窗口说明采用订单日期、交付日期还是快照日期；
  \item 对空值、未交付订单和分母为零的情况给出处理规则；
  \item 查询输出字段命名稳定，能够进入后续看板或报告；
  \item 至少提供一个行数检查或 Python 复核结果。
\end{enumerate}

上述检查同时提出了对输入数据的要求。重复键改变连接次数，缺失外键改变记录覆盖范围，日期错误改变统计窗口。第~\ref{ch:week04-data-quality}~章研究这些问题的生成机制及其对指标的影响。

\section*{本章小结}
\addcontentsline{toc}{section}{本章小结}

SQL 通过选择、连接和聚合把业务记录转化为统计量。连接键的唯一性与外键覆盖决定记录是否被复制或遗漏，分组字段决定结果粒度，空值处理决定进入计算的对象范围。窗口函数在保留记录身份的同时计算局部或累计统计量。

聚合状态给出了一个可检验的原则：先保留足够的可合并信息，再在目标粒度上计算指标。制造场景中，良率需要计数权重，OEE 需要时间与理想节拍的一致口径，WIP 需要区分快照和时间平均，OTD 需要声明到期还是已交付队列。AI 问数必须继承这些约束，并由独立执行层落实授权和资源限制；标准查询与自动测试用于核对口径、数据范围和计算结果。比较不同平台时，还要固定问题集、数据快照、用户角色和时限。

查询正确性以输入记录满足相应约束为条件。记录缺失、测量误差与规则违反均可能改变统计结论，需要结合数据生成过程分析。

\chapterreading{Chamberlin 与 Boyce 的 SEQUEL 论文有助于理解声明式查询的历史动机
\cite{chamberlin1974}。关系语义与数据库基础见 Date\cite{date2003}，常见设计
反模式可参考 Karwin\cite{karwin2010}。制造指标应结合 SEMI 的设备绩效导读及
ISO 22400-2 的定义范围阅读\cite{semi03,iso22400}，并确认所采用的时间分类与设备范围。Spider 2.0 与 BIRD-INTERACT 研究多步任务及交互查询的评价，评价对象超出
单条 SQL 的文本匹配\cite{spider03,bird03}。平台语义定义与问数接口的实现可对照
\cite{databricksGenie2026,snowflakeSemantic2026,snowflakeAgents2026}。}

\section*{思考与习题}
\addcontentsline{toc}{section}{思考与习题}

\begin{exercise}[制造指标反例]
构造两个不同理想节拍的产品，使按件数汇总的 \(APQ\) 与理想良品工时比不同；说明差异产生在哪个权重。再构造每天两个产品族的 WIP 快照，比较行均值与日总量均值。
\end{exercise}
\begin{exercise}[问数评测]
为“上周 OTD 最低的工站”设计五个测试：同义改写、零分母、逾期未交付、重复连接和越权范围。写明每个测试的期望行为，区分需要澄清、拒绝执行和可回答三类结果。
\end{exercise}

\begin{exercise}
使用 \path{data/sample/orders.csv}、\path{data/sample/customers.csv} 和 \path{data/sample/products.csv}，按客户分层计算销售额、毛利率和准时交付率。要求写出 SQL，并说明事实表粒度、分组粒度、分子、分母和时间窗口。
\end{exercise}

\begin{exercise}
解释为什么不能直接平均各区域毛利率来得到总体毛利率。请构造一个两区域的小例子，给出错误算法和正确算法的数值差异。
\end{exercise}

\begin{exercise}
写出一个 SQL 查询，找出每个客户最近三笔订单的累计收入。请说明 \texttt{PARTITION BY}、\texttt{ORDER BY} 和窗口框架分别承担什么作用。
\end{exercise}

\begin{exercise}
对 \texttt{orders} 与 \texttt{products} 的连接结果做行数检查。若连接后行数增加，说明可能出现了什么问题？若连接后行数减少，又可能是什么原因？
\end{exercise}

\begin{exercise}
为“高优先级订单延期率”设计一条指标字典记录，至少包含对象范围、事实粒度、分子、分母、时间窗口、单位、维度和责任部门。
\end{exercise}
